\section{A fixed rational field and its finite certificates}\label{app:repro}

\begin{proposition}[Six cycles at specified coefficients]\label{prop:fixed-six}
In~\eqref{eq:deg12-poly} set $d=6/5$, $\sqrt\eta=10^{-40}$,
$\theta=-10^{-20}$ and $a=(a_-+a_+)/2$ with $a_\pm$ from~\eqref{eq:seed-coeffs}.
The resulting rational primitive has degree twelve. Its field has a
hyperbolic cycle in each bracket of Table~\ref{tab:fixed-six} and a
repelling hyperbolic cycle above those five.
\end{proposition}

For a forward arc of the reflected field with $G_s(x)=F(sx)$, let
$\Lambda_s(h)=\int G_s'(x(\tau))\,\dd\tau$ from launch $(0,h)$ to
its first lower-axis hit $(0,-w_s(h))$. This is the negative of the
integrated divergence on that half-arc.

\begin{table}[htbp]
\centering\small
\caption{Physical upper-axis brackets for the fixed negative-tilt field.
Every decimal endpoint is rational. The last column encloses the
integral difference over the whole closed bracket.}\label{tab:fixed-six}
\begin{tabular}{@{}rrrrc@{}}
\toprule
$h_l$ & $h_r$ & $\sgn\D(h_l)$ & $\sgn\D(h_r)$ & $10^5(\Lambda_+-\Lambda_-)$\\
\midrule
3.494474846805 & 3.494474846808 & $+$ & $-$ & $(235,236)$\\
3.848603188665 & 3.848603188668 & $-$ & $+$ & $(-4.1,-3.6)$\\
3.978624148391 & 3.978624148394 & $+$ & $-$ & $(1.5,2.1)$\\
4.324131925377 & 4.324131925381 & $-$ & $+$ & $(-1.5,-0.9)$\\
4.454460589336 & 4.454460589341 & $+$ & $-$ & $(1.8,2.6)$\\
\bottomrule
\end{tabular}
\end{table}


\begin{proof}
Liouville's formula with the normal velocities at the two sections
gives $w_s'(h)=(h/w_s(h))e^{-\Lambda_s(h)}$. At a zero of
$\D=w_+-w_-$ the full multiplier is therefore
$e^{-(\Lambda_+-\Lambda_-)}$.

\emph{Compact brackets.}
Let $Q_s$ be the sextic from~\eqref{eq:seed-coeffs} at the same midpoint,
reflected on side $s$. Summing the absolute differences of the
coefficients weighted by $3^j$, first for $G_s-Q_s$ and then for its
derivative, gives a bound $\delta=2\cdot10^{-74}$ for each on $|x|\le3$.
The leading coefficient of the full primitive is negative.
For each rational endpoint and for each entire height bracket we use
the order-36 dyadic paths of Lemma~\ref{lem:seed}. In normalized time
$0\le v\le1$ let $R$ bound the residual against the sextic vector
field multiplied by the physical step $b$ (the residual is the path
derivative minus $b$ times that field). Its polynomial coefficients
through degree 216 are formed exactly on rational data. Round path
coefficients to the $2^{-256}$ grid and error bounds outwards.
The symmetric part of the sextic Jacobian is
$\operatorname{diag}(-Q_s',0)$. The completed squares of
Lemma~\ref{lem:seed} give logarithmic norm at most $\gamma=1$ for
$x\ge-1/100$, and at most $1/250$ for $x\ge0$. Thus a valid radius is
\[
 e_{\rm out}=\frac{e_{\rm in}+R+b\delta}{1-b\gamma}.
\]
Indeed $e'\le b\gamma e+R+b\delta$ until first exit.
Each accepted tube lies inside its domain; the smaller $\gamma$ is
used only after verification of positivity with the wider tube.
Centering the initial tube at the bracket midpoint with radius half
its width covers every launch therein.

Compose $Q_s'$ with the proposed $x(v)$, integrate its polynomial
on each step, and enclose the coefficients by outward dyadic rounding.
The additional radius in the integral is at most
$b(\|Q_s''\|e_{\rm out}+\delta)$. The departure and intermediate
tubes have $x>0$ after the initial crossing. On the final step the
enclosed arc crosses $x=0$ with $-\dot x\ge v_0>0$. At the proposed
event time found by bisection the true time differs by at most
$(|x_{\rm app}|+e_{\rm out})/v_0$ in physical units. Multiply this
bound by $\sup|G_s'|$ for the last integral and by $\sup|x|$ for the
hit ordinate, in addition to the path error there. The resulting endpoint
signs and full-interval integral differences are in Table~\ref{tab:fixed-six}.
Order and continuity give a zero in each interval and the displayed
multiplier formula proves simplicity.

\emph{Outer section hits at the fixed scale.}
Set $P_s(z)=F(sLz)/S$, $u=z^2/2$, $p_s(u)=P_s(\sqrt{2u})$,
$t=y/S$, $\eps=L^2/S^2$, and put
$\lambda=\sqrt\eta$, $K=10^7$, $c=10^{-4}$, $\rho=\lambda^4$, $\nu=\rho^2$.
The rational profile bounds needed below are as follows. The maximum
$M_s$ on $0\le u\le2$ occurs at $1.18<z_m<1.22$,
$0.012<M_s<0.014$, and
\[
 P_s(d)\le M_s\le P_s(d)+m_{\rm lift},\qquad
 m_{\rm lift}=(10^5\eta)^2/(2\cdot0.1).
\]
On $1.18\le z\le1.22$ one has $p_s''<-0.1$ with opposite signs of
$p_s'$ at the ends; on $1.22\le z\le2$ one has $p_s'<-0.01$,
$P_s(1.22)>0.01$, $P_s(2)<-2$. Further,
$P_s>10^{-3}$ on $[1/2,1.22]$, $P_s<M_s-2c$ on $[0,1.11]$,
$|p_s'|,|p_s''|<K$ on $1/2\le z\le2$, and
$\int_0^{1.10}|P_s'|\,\dd z<2K$. These are finite polynomial
comparisons. Use $p_s'=P_s'/z$, $p_s''=(zP_s''-P_s')/z^3$ away
from zero. For a polynomial $Q(a+(b-a)v)=\sum_j c_j v^j$ of
degree $k$, its Bernstein coefficients
$\sum_{j\le i}\binom{i}{j}c_j/\binom{k}{j}$ bound it on $0\le v\le1$.
The explicit even profile increases up to $d$; coefficient norms
bound its perturbation, yielding the stated height margins below
$z=1.18$. At $d$, $|p_s'(d^2/2)|\le10^5\eta$ and the curvature
bound yields the lift. Translating $P_+(z)+10^{-3}z^{12}$ by $2$
gives negative coefficients, proving the tail bound for $z\ge2$.

On the plus side follow a launch from $(z,t)=(1/2,y/S)$,
$-1\le y\le9$, backward in physical time. Then $t$ increases and
$\eps u_t=p_+(u)-t$. The falling root $U(t)$ satisfies $|U'|\le1/\rho$
up to $t=M_+-\rho$. Indeed its distance in $u$ from the maximum exceeds
$20\rho$ there, since $200K\rho^2<\rho$. The curvature and outer
slope bounds give $-p_+'\ge\rho$ on $[U(t)-\nu,2]$.
The tube $|u-U(t)|\le\nu$ is invariant. Its boundary speed toward the
root is at least $\rho\nu/\eps$. At launch the orbit is left of
the tube. The bounds through $z=1.22$ and monotonicity thereafter
give at least this speed until capture. Even allowing for the root
speed, it covers distance $2$ within a height increase $\rho$ since
$\eps<\rho^4/8$; the walls at $z=1/2$ and $u=2$ confine this
initial passage. At $M_+-\rho$ the orbit is still away from the axis.
From $t=M_++\rho$ onward it moves inward at speed at least
$\rho/\eps$ and reaches the axis within $2\eps/\rho<\rho$
if it has not already crossed. Boundedness and positive height near
exit exclude an asymptotic stop. Write $\widehat h_+(y/S)$ for this
scaled upper hit. Near exit, from $z=1.10$ to $0$, the inverse equation
$\dd t/\dd u=-\eps/(t-p_+(u))$ has denominator $>c$.

Order gives two physical heights
$H_1=S\widehat h_+(9/S)<H_2=S\widehat h_+(-1/S)$: the earlier
launch lies to the right of the section when the later one starts.
Both are strictly inside the rational enclosure
\[
 \mathcal I=\big(S(P_+(d)-\rho),\ S(P_+(d)+m_{\rm lift}+2\rho)\big).
\]
Its lower endpoint exceeds $4.6$. At the rational upper endpoint
$t\ge M_++2\rho$ initially. In forward time $\eps\,\dd u/\dd(-t)=t-p_+$,
so the loss in height to $u=2$ is at most $2\eps/\rho<\rho$.
From there the tail bound gives at most
$\eps\int_2^\infty[10^{-3}(2u)^6]^{-1}\dd u=25\eps/256<\rho$
before escape at positive height. The sign tests use $H_1,H_2$.

\emph{Signs at the two hits.}
The rational comparison $SP_-(d)>\sup\mathcal I$ makes $x=Ld$ a
forward wall on the minus side. On both sides $G_s(x)>8x^6$ for
$2\le x\le L/2$ and $G_s(x)>x^6/100$ on $L/2\le x\le Ld$.
For these tests use the sextic coefficient error on $[2,10]$ with
weights $10^j$, the factorization of $E'$ in~\eqref{eq:deg12-even}
from $10$ to $L/2$, and $P_s>10^{-3}$ thereafter.
At $8\le y\le9$, $v=x^2/2$ obeys $\dd v/\dd(-y)=y-G_s(x)$.
For $v\ge50$ this is at most $-v^3/25$; for $2\le v\le50$ it is
at most $-4v^3$. The descent needed to reach $x\le2$ is less than
$25/(2\cdot50^2)+1/32<1$. The boundary at $x=2$ points inward
throughout this range. Thus at $y=8$ the minus arc, as well as the
plus arc through $(L/2,9)$, lies in $0\le x\le2$. The wall at $x=3$
and the focus condition $|a|<2$ give compact returns. For their hit
magnitudes $B_s(q,8)$ from $(q,8)$ the compact tubes above give
\[
 B_+(2,8)-B_-(0,8)<-2\cdot10^{-32},\qquad
 0.03<B_-(0,8)\le B_-(2,8)<0.1.
\]
Monotonicity in $q$ yields $\D(H_1)<0$. From $(L/2,-1)$ the plus
arc moves inward to $x=2$ by the same tail bounds and returns inside
the compact wall with magnitude $>1$. Hence $\D(H_2)>0$. Order
and the bounding returns give continuity throughout $[H_1,H_2]$.
Any zero in the internal bracket has common lower magnitude $<0.1$
and plus section ordinate $y\in(-1,9)$.

\emph{Multiplier of the outer zero.}
In scaled coordinates $\eps \Lambda_s=\int p_s'(u)\,|\dd t|$.
On the backward-tracked part of the plus arc the slope is negative;
for $t\in[0.002,0.003]$ it is below $-0.01$ since the tube lies
beyond $z=1.22$. This contributes at most $-10^{-5}$. The entrance
and exit layers above $z=1.10$ cost at most $4K\rho$ in total, and
the final inverse chart costs at most $2K\eps/c$ by the bound
on $\int |P_+'|$. On the forward leg from $L/2$ to $x=2$ one has
$|p_+'|<K$: on $0.01/L\le z\le d$ bound the singular linear
coefficient divided by $z$ using $z\ge0.01/L$, and bound the
remaining polynomial by its coefficient norm. The available physical
height change is less than $9.1$, giving $10K/S$.

On the minus side the initial passage to $x=0.01$ has $t-p_->0.009$
while $t\ge0.01$; the height loss is at most
$\eps(0.01/L)^2/(2\cdot0.009)<0.002$. The inverse chart contributes
at most $2K\eps/0.009$ in absolute value. The boundary at $x=0.01$
points right down to $y=8$. There $z\le d$ by the wall. On this
portion $E'(z)/z\ge0$ and the correction to the slope has magnitude
$<10^{-8}$ using $z\ge0.01/L$. The possible negative contribution
is at most $0.014\cdot10^{-8}$ in magnitude.

For either remaining compact arc, $|G_s'|<10^5$ on $[0,3]$.
Above $x=0.01$ the time is at most $9.1/0.01$. After crossing
inward through $x=0.01$ the monotonicity of $y$ prevents recrossing.
Here $|G_s(x)/x|<0.2$ gives angular velocity at most $-0.9$, so
the remaining time is less than five. Each compact part contributes
less than $10^8\eps$ in absolute value. Substitution of the fixed
scales gives
\begin{align*}
 \eps \Lambda_+ &\le -10^{-5}+4K\rho+2K\eps/c+10K/S+10^8\eps<-9\cdot10^{-6},\\
 \eps \Lambda_- &\ge -2K\eps/0.009-0.014\cdot10^{-8}-10^8\eps>-10^{-6}.
\end{align*}
Thus $\eps(\Lambda_+-\Lambda_-)<-8\cdot10^{-6}$ at every zero in
$(H_1,H_2)$. Its multiplier exceeds one. This bracket and the five
compact brackets are disjoint.
\end{proof}

\subsection*{What the fixed-scale matching tests resolve}
At $\lambda=\sqrt\eta=10^{-40}$ the signs on the parameter faces give
a real solution of the three matching equations. At a match the entrance comparisons
at $4/5,0,-20$ have signs $-,+,-$ and the direct launch gives $+$.
The bounds below give necessary tolerances for mixed shoulder signs
and for paired section hits in the stated windows. The eight shoulder
signs used for nineteen cycles follow from the moving-fold estimate
as $\eta\to0$ in Proposition~\ref{prop:moving-layer}. For the fixed-scale
estimates take the closed parameter box
\[
 1.19888999543\le d\le1.19888999544,\qquad
 |\theta|\le10\lambda,\qquad a_-\le a\le a_+ .
\]
With $\beta=-K/\widehat B'(1)$ as in Section~\ref{subsec:deg12-prep},
the scales $L=-a_5/(6a_6\beta\eta)$, $S=6a_6L^6$ and $\eps=L^2/S^2$
from~\eqref{eq:deg12-poly} give $4.07\cdot10^{-781}<\eps<4.09\cdot10^{-781}$
throughout this box.
Write $H_s^*$ for the scaled upper hit traced backwards from
$(u,t)=(0,-1)$. Let $g_s(r)$ be its positive separation from the
upper hit traced from $(1/8,r)$, and $g_s^b$ its separation from
the hit traced from $(0,-1/2)$.

\emph{The common outer height.}
On either side let $U(t)$ be the falling root of $p_s(U)=t$, let
$z_m$ denote the maximizer and put $v_s(r)=\sqrt{2U(r)}$.
The curvature, wall and tube bounds used in
Proposition~\ref{prop:fixed-six} hold on the whole box by the same
Bernstein and coefficient inequalities. For $-22\le Sr\le1$,
and also for the direct launch with $r=-1/2$, they give
\begin{equation}\label{eq:finite-gap-cost}
 \left|-\eps\log g-
   \int_{z_m}^{v_s(r)}\frac{P_s'(z)^2}{z}\,\dd z\right|<10K\rho .
\end{equation}
Indeed the reference has already reached its tube by the second launch.
The initial $u$-separation lies in $(1/2,2)$. For the positive
separation $b(t)$, $\eps(\log b)'$ equals the secant of $p_s$.
The launch layer costs at most $K\rho$ in the secant integral.
On the common tube the secant differs from $p_s'(U(t))$ by at most
$K\nu$. Omitting the first and last $\rho$ of root tracking costs at most
$2K\rho$. Up to the first inward crossing of $u_c=(1.10)^2/2$ the
remaining layer costs at most $3K\rho$. The secant is bounded there
because at least one $z$ exceeds $1/2$, so $2/(z_1+z_2)\le4$.
Contraction on $[0.002,0.003]$ makes the separation $b_c$ at that
crossing smaller than $2e^{-5\cdot10^{-6}/\eps}$. Both trajectories
then lie below $z=1.11$, with inward speed between $c/\eps$ and
$K/\eps$. Their crossing heights differ by between
$\eps b_c/K$ and $\eps b_c/c$. In the last inverse chart,
the logarithmic derivative of their height difference with respect to $u$
is $\eps/[(t_1-p_s)(t_2-p_s)]\le\eps/c^2$. In $\eps\log$
units this conversion and the initial factor together cost less than
$6000\eps<K\rho$. Finally $t=P_s(z)$ gives the integral in
\eqref{eq:finite-gap-cost}.

For a rational enclosure, bracket $z_m$ by the signs of $P_s'$ on
$[1.18,1.22]$ and the falling root by those of $P_s-r$ on $[1.22,2]$.
Integrate the polynomial part of $P_s'^2/z$ exactly. The remaining
term is $P_s'(0)^2\log(v_s/z_m)$ with $0<\log(v_s/z_m)<1$.
Including the error in~\eqref{eq:finite-gap-cost} gives
$0.0019<-\eps\log g_s(r)<0.0021$ and
$1.3<-\eps\log g_s^b<1.4$. Thus equality of a pair of section
hits in this window forces
\[
 |H_+^*-H_-^*|=|g_+(r_+)-g_-(r_-)|<e^{-4\cdot10^{777}} .
\]
The direct hit is within $e^{-3\cdot10^{780}}$ of its own reference.

\emph{The inner shooting parameter.}
Put $T=(E(1)+E(d))/2$ and fix $d,\theta$. At $a_-,a_+$ respectively
the compact enclosures at $y=8$ give
$\D(ST)>2\cdot10^{-32}$ and $\D(ST)<-6\cdot10^{-32}$, including the
uniform $10^{-72}$ perturbation of the sextic on $[0,3]$ in the
Taylor defect. The linear-parameter sensitivity in
Lemma~\ref{lem:rotational} gives a unique $a_*$ with $\D(ST)=0$.
Consider a shoulder launch from the eight levels in
Proposition~\ref{prop:moving-layer} and the launch at $T$ on the same side.
The shoulder levels lie in $(0.0108,0.0110)$ below $T$.
The wall $z=1.18$ confines both launches. For $0.006\le t\le0.007$,
outside $1/2\le z\le0.9$ the forward equation
$\eps\,\dd u/\dd(-t)=t-p_s(u)$ points toward that strip with magnitude
at least $10^{-4}$ on the right-hand side. Indeed $P_s<0.0019$ below its
left edge and $P_s>0.0071$ above its right edge up to the wall, by
the even profile and its perturbation bound. The distance in $u$ is
less than $2$, so both arcs are inside by $t=0.006$. The boundary
signs persist down to $0.002$, and on the strip $p_s'>0.005$.
On the remaining descent to $y=8$, $z\le0.9$ and $x\ge0.1$ by the
walls. The even slope is nonnegative and the correction exceeds
$-10^{-8}$ there (use $z\ge0.1/L$ in the linear term).
Integrating the secant equation in decreasing $t$ gives a separation
in physical $x^2/2$ below $2L^2e^{-1.9\cdot10^{-5}/\eps}$ at $y=8$.

On this last section the map from $x^2/2$ to the hit magnitude has
derivative $e^{-\Lambda_s}/w_s$ by Liouville's formula, with the
integral and the terminal magnitude now taken from $(x,8)$. For $0\le x\le2$,
order and the endpoint tubes give $0.03<w_s<0.1$. The $x=0$ launch
returns within $41$ time units, allowing $10^{-29}$ for the variation
of the full field over the linear-parameter interval. A larger arc
reaches that terminal ordinate at least as fast, since its $x(y)$ is
larger. Afterwards $\dot x\le-0.03+0.015$ by $G_s'>-0.005$ on
$[0,3]$, so the remaining time is at most $200$. The derivative is
therefore at most $e^{1.25}/0.03<200$. Each shoulder hit lies below
the corresponding reference in magnitude, by order. The two losses
are each less than $400L^2e^{-1.9\cdot10^{-5}/\eps}
<e^{-4\cdot10^{775}}$. Opposite signs at two shoulder heights force
$|\D(ST)|$ below that same bound.

In the formula for $|\partial_a w_s(ST)|$, retain the first
$5\cdot10^{-7}$ time units after $y=8$. There $x\ge0.05$ by
$|\dot x|<10^5$. The remaining compact time is less than $250$ and
$|G_s'|<10^5$. Thus $-\partial_a\D(ST)>e^{-3\cdot10^7}$ throughout
the shooting interval. The mean value theorem gives the necessary
condition
\[
    |a-a_*(d,\theta)|<e^{-3\cdot10^{775}}
\]
for mixed shoulder signs at these fixed $d,\theta$. These estimates
concern the specified box and these choices of sections.

The remaining tests concern the faces of this box and the signs of
individual outer samples. For example, at the two $d$ faces,
\eqref{eq:finite-gap-cost} bounds $(\eps/\eta)\log[g_+(\Delta/S)/g_-(0)]$
in $(-2.974,-2.972)10^{-13}$ and $(4.704,4.707)10^{-13}$, respectively.
At the two $\theta$ faces the peak-lift bound places
$(H_+^*-H_-^*)/\lambda^3$ in $(-46.75,-46.74)$ and $(92.96,92.97)$.
Together with the shooting faces these signs give a real match by the
same box argument as in Section~S5.

After matching, section ordinates require finer control. At $b=0.003$
write $b_s(r)$ for the separation in $u$ from the reference. A section
launch reaches $z=1.22$ within $1500\eps$ in height since
$p_s-t>0.0009$ on the launch segment. Thereafter
$-p_s'\ge0.01$ between that wall and $z=2$, and $|U'|\le100$.
Variation of constants gives distance at most $10002\eps$ from $U$
after $400000\eps$ in height, including the launch layer.
The secant integral then gives
$|\eps\log b_s(r)-\int_r^b p_s'(U(t))\,\dd t|<10^{13}\eps$:
the layer costs at most $2K(400000\eps)$ and the tracked error at
most $10002K\eps$. From $b$ to the inward crossing of $u_c$ the
whole interval stays within width $d_*=2e^{-2\cdot10^{-6}/\eps}$,
using the contraction to $b$ and expansion at most $3K\rho$ in action
thereafter. Its crossing times differ by at most $\eps d_*/c$.
For the map from the section at $b$ to the scaled upper height $H$,
the variational equation gives its derivative as
$(\eps/H)\exp(\eps^{-1}\int_b^H p_s'(u(t))\,\dd t)$.
One can first stop at positive $u$ and then take the limit in the
inverse chart, because $\int_0^{u_c}|p_s'(u)|\,\dd u<2K$ there.
In the common time interval the action difference is at most $Kd_*$,
and in the remaining time interval it is at most $K\eps d_*/c$.
In the inverse chart the height difference contracts. Substituting
$\dd t=-\eps\,\dd u/(t-p_s)$ bounds the difference of the final
actions by $2K\eps^2d_*/c^3$. The transverse prefactor in the
derivative of the upper hit is $\eps/H$, with $H>0.01$, so its
contribution to the oscillation of the scaled log derivative is at most
$100\eps^2d_*/c$. Together these contributions are below $5Kd_*$.
Taking the ratio for two launches now cancels the common terminal
derivative up to this oscillation. With $E(v_0)=0$ on the falling
branch and $c_d=-E'(v_0)/v_0>0.01$, it gives
\[
 \left|S\eps\log\frac{g_s(y_1/S)}{g_s(y_0/S)}
          -c_d(y_1-y_0)\right|<10^{-55},\qquad -25\le y_i\le1 .
\]
Indeed the error from the preceding estimates is below
$3\cdot10^{13}S\eps<10^{-280}$. The slope error over the physical
interval of length $26$ is bounded by
$26(\|R_s'\|+100K(\|R_s\|+25/S))<10^{-55}/2$, with $R_s=P_s-E$
and coefficient norms on $[0,2]$. At a match, tests at
$y-\Delta\pm10^{-12}$ therefore pair each of $4/5,0,-20$ with an
opposite hit. The compact tubes with the $1/448$ tail allowance give
their signs $-,+,-$. For the direct launch one uses the first variation
of the action in~\eqref{eq:finite-gap-cost} along $E+\gamma R_s$.
The peak term vanishes and the falling endpoint supplies
$-(P'/z)R_s$ there. With $\mu=10^7\eta$ bounding $R_s,R_s'$,
the falling root moves at speed at most $100\mu$ in $\gamma$.
Differentiating the integral a second time bounds its interior term
by $2\mu^2$ and the moving-boundary and endpoint terms together by
$(1+500K)\mu^2$. The two Taylor remainders thus cost less than
$10^5K\mu^2$. Including $20K\rho$ from the two gap estimates in the
first-order comparison of Section~S5 gives
$3.8<(\eps/\eta)\log(g_+^b/g_-^b)<5.5$. With common references
the plus direct hit is lower, so the minus direct height has positive
return sign. These tests certify particular comparisons at the fixed
scale. The eight moving-fold signs used together in
Theorem~\ref{thm:deg12} come from the asymptotic estimate of
Proposition~\ref{prop:moving-layer}.

At the two specified midpoints with $d=6/5$ and
$\theta=\pm10^{-20}$, outside the box, the maxima differ already in
height. More precisely $H_+^*-H_-^*$ is in $(6.99,7)10^{-100}$ for
positive tilt and in $(-7,-6.99)10^{-100}$ for negative tilt. These
intervals follow from $P_s(d)\le M_s\le P_s(d)+m_{\rm lift}$ and
$M_s-\rho<H_s^*<M_s+2\rho$. The six-cycle certificate uses the two
interior hits of Proposition~\ref{prop:fixed-six}. The nineteen-cycle
argument uses the common matches and regularization of Section~S5.

\paragraph{Reproducing the finite inequalities.}
Unpack the certificate archive and verify its contents before running the programs.
\begin{verbatim}
tar -xzf certificates.tar.gz
cd certificates
sha256sum -c CERTIFICATE-SHA256SUMS
./replay-all.sh
python3 -B verification/check_action_coordinates.py
\end{verbatim}
The programs use Python's standard library. In \texttt{REPLAY.txt},
\texttt{verify\_paper\_tables.py} checks the printed cubic and well
root brackets, and \texttt{verify\_one\_sided\_wells.py} checks the
octic moments. For Theorem~\ref{thm:deg12}, \texttt{verify\_seed.py}
encloses all fourteen sextic half-arcs in Table~\ref{tab:seed};
\texttt{verify\_entrance.py} bounds the tail from infinity to $x=2$
by $1/448$. The joint-shoulder program checks the retained actions and
turning comparison. The three matching equations, choice of regular
value and nineteen zeros are proved in Section~\ref{subsec:deg12-finish}.
For Proposition~\ref{prop:fixed-six} the directory
\nolinkurl{verification/fixed-degree12/} contains the thirteen rational
coefficients in \nolinkurl{inputs/outer-variant-field.json}, the five
height brackets, the polynomial checks and the validated integrator.
The fixed-field replay prints each integral enclosure and the outer
inequality in \nolinkurl{verification/fixed-degree12-replay.txt}.
Changing only $\theta$ in Proposition~\ref{prop:fixed-six} to
$+10^{-20}$ gives the positive-tilt field, for which the same five
compact brackets are certified. The weak-damping field of
Appendix~\ref{sec:weak} has its own coefficients and brackets.
